\documentclass[11pt,letterpaper]{article}
\usepackage[margin=1in]{geometry}
\usepackage{amsmath,amssymb}
\usepackage{graphicx}
\usepackage{booktabs}
\usepackage{hyperref}
\usepackage{fancyhdr}
\usepackage{titlesec}
\usepackage{float}

% Header and footer setup
\pagestyle{fancy}
\fancyhf{}
\lhead{\textbf{PROJECT_NAME} - Analysis Report}
\rhead{Page \thepage}
\renewcommand{\headrulewidth}{0.4pt}

% Title formatting
\titleformat{\section}{\Large\bfseries}{\thesection}{1em}{}
\titleformat{\subsection}{\large\bfseries}{\thesubsection}{1em}{}

% Document metadata
\title{\textbf{PROJECT_NAME}\\Analysis Report}
\author{Arcanum Research Initiative}
\date{\today}

\begin{document}

\maketitle

\tableofcontents
\newpage

\section{Executive Summary}

\textbf{Analysis Period:} [Time period covered]

\textbf{Key Findings:} [Primary results in bullet points]
\begin{itemize}
    \item [Finding 1 with quantified result]
    \item [Finding 2 with quantified result]
    \item [Finding 3 with quantified result]
\end{itemize}

\textbf{Policy Implications:} [Practical applications of results]

\textbf{Statistical Confidence:} [Accuracy metrics and validation results]

\section{Data Overview}

\subsection{Dataset Summary}
\begin{table}[H]
\centering
\begin{tabular}{@{}lcc@{}}
\toprule
\textbf{Data Source} & \textbf{Observations} & \textbf{Time Period} \\
\midrule
[Source 1] & [Number] & [Period] \\
[Source 2] & [Number] & [Period] \\
[Source 3] & [Number] & [Period] \\
\bottomrule
\end{tabular}
\caption{Data Sources and Coverage}
\end{table}

\subsection{Data Quality Assessment}
\begin{itemize}
    \item \textbf{Completeness:} [Percentage of complete records]
    \item \textbf{Accuracy:} [Validation results]
    \item \textbf{Consistency:} [Cross-source validation]
\end{itemize}

\section{Primary Analysis}

\subsection{[Analysis Topic 1]}

\textbf{Objective:} [What this analysis aims to determine]

\textbf{Methodology:} [Brief approach description]

\textbf{Results:} [Quantified findings]

% Include visualization
\begin{figure}[H]
    \centering
    \includegraphics[width=0.8\textwidth]{path/to/figure1.png}
    \caption{[Descriptive caption]}
    \label{fig:analysis1}
\end{figure}

\textbf{Interpretation:} [What these results mean]

\subsection{[Analysis Topic 2]}

\textbf{Objective:} [What this analysis aims to determine]

\textbf{Methodology:} [Brief approach description]

\textbf{Results:} [Quantified findings]

% Include results table
\begin{table}[H]
\centering
\begin{tabular}{@{}lcc@{}}
\toprule
\textbf{Category} & \textbf{Value} & \textbf{Confidence} \\
\midrule
[Category 1] & [Value] & [Confidence level] \\
[Category 2] & [Value] & [Confidence level] \\
[Category 3] & [Value] & [Confidence level] \\
\bottomrule
\end{tabular}
\caption{[Analysis Results Summary]}
\end{table}

\textbf{Interpretation:} [What these results mean]

\section{Comparative Analysis}

\subsection{Benchmarking Results}

\textbf{Comparison Framework:} [What we're comparing against]

\textbf{Performance Metrics:}
\begin{itemize}
    \item [Metric 1]: [Result vs benchmark]
    \item [Metric 2]: [Result vs benchmark]
    \item [Metric 3]: [Result vs benchmark]
\end{itemize}

\begin{figure}[H]
    \centering
    \includegraphics[width=0.8\textwidth]{path/to/comparison.png}
    \caption{Comparative Analysis Results}
    \label{fig:comparison}
\end{figure}

\section{Statistical Validation}

\subsection{Accuracy Assessment}
\begin{itemize}
    \item \textbf{Cross-Validation Results:} [Accuracy percentage]
    \item \textbf{Error Analysis:} [Types and magnitude of errors]
    \item \textbf{Confidence Intervals:} [Statistical confidence ranges]
\end{itemize}

\subsection{Robustness Testing}
\begin{itemize}
    \item \textbf{Sensitivity Analysis:} [How results change with assumptions]
    \item \textbf{Alternative Methodologies:} [Comparison with other approaches]
    \item \textbf{Outlier Assessment:} [Treatment of extreme values]
\end{itemize}

\section{Detailed Findings}

\subsection{[Finding Category 1]}

\textbf{Description:} [Detailed explanation of finding]

\textbf{Evidence:} [Supporting data and analysis]

\textbf{Implications:} [What this means for the research question]

\subsection{[Finding Category 2]}

\textbf{Description:} [Detailed explanation of finding]

\textbf{Evidence:} [Supporting data and analysis]

\textbf{Implications:} [What this means for the research question]

\section{Policy and Practical Applications}

\subsection{Key Insights for Stakeholders}
\begin{itemize}
    \item [Actionable insight 1]
    \item [Actionable insight 2]
    \item [Actionable insight 3]
\end{itemize}

\subsection{Recommended Actions}
\begin{enumerate}
    \item [Specific recommendation 1]
    \item [Specific recommendation 2]
    \item [Specific recommendation 3]
\end{enumerate}

\section{Limitations and Caveats}

\subsection{Data Limitations}
\begin{itemize}
    \item [Limitation 1 and impact]
    \item [Limitation 2 and impact]
    \item [Limitation 3 and impact]
\end{itemize}

\subsection{Methodological Considerations}
[Discussion of analytical constraints and assumptions]

\section{Future Research Directions}

\subsection{Extension Opportunities}
\begin{itemize}
    \item [Research extension 1]
    \item [Research extension 2]
    \item [Research extension 3]
\end{itemize}

\subsection{Data Enhancement Priorities}
[Areas where additional data would improve analysis]

\section{Conclusion}

\textbf{Primary Conclusions:} [Main takeaways from analysis]

\textbf{Research Contribution:} [How this adds to existing knowledge]

\textbf{Practical Value:} [Real-world applications and benefits]

\section{Data and Code Availability}

\textbf{Dataset Files:}
\begin{itemize}
    \item \texttt{Output/Data/[filename].xlsx} - [Description]
    \item \texttt{Output/Data/[filename].xlsx} - [Description]
\end{itemize}

\textbf{Analysis Code:} Available in \texttt{Technical/src/} directory

\textbf{Reproduction Instructions:} See Technical README for step-by-step guide

\section{Appendix}

\subsection{Additional Tables}
[Supplementary data tables]

\subsection{Technical Details}
[Detailed methodological information]

\subsection{Data Dictionary}
[Complete variable descriptions for all datasets]

\end{document}